%%%%%%%%%%%%%%%%%%%%%%%%%%%%%%%%%
% Conclusion
%%%%%%%%%%%%%%%%%%%%%%%%%%%%%%%%%

\section{Conclusion}
\label{sec:conclusion}

Two axioms, element-wise independence and conditional independence of block error rates, yield the distributional results collected here: binomial error counts, normal limits, and probability intervals for the realized rates.
These are standard consequences of independence, stated once so that the companion papers can cite them.
For a single second-order Bernoulli set, the Kronecker factorization of the joint confusion matrix means that the probabilistic structure is determined by the false positive and false negative rates.
For composed set expressions it is not, because the rates depend on each element's latent membership pattern (\cref{sec:higher_order_model}).
Formulating the model as an abstract data type means that any implementation satisfying the axioms inherits these results; the remark ``\nameref{rem:implementations}'' in \cref{sec:bernoulli_model} describes the conditional sense in which a Bloom filter does so.

\paragraph{Practical implications.}
The abstract data type formulation means that the results of this paper apply automatically to any implementation whose membership queries satisfy element-wise independence.
This separation of the probabilistic guarantee from the implementation details is the key advantage of the axiomatic approach.

\paragraph{Open problems.}
The model assumes element-wise independence of errors.
Extending the framework to \emph{correlated} error models---such as those arising in locality-sensitive hashing or cache-line-aligned filters---would broaden applicability, though the clean binomial structure would likely give way to more complex combinatorics.

\paragraph{Companion papers.}
The set-theoretic composition of Bernoulli sets---closed-form error rates for complement, union, intersection, and set difference, monoidal structure, relational predicates, and interval arithmetic---is developed in~\cite{bernoulliComposition}.
The probability distributions of binary classification measures---positive predictive value, negative predictive value, accuracy, and Youden's $J$---induced by the Bernoulli model are derived in~\cite{bernoulliMeasures}.
The joint entropy of error counts and the information-theoretic space complexity lower bound for approximate sets are developed in~\cite{bernoulliEntropy}.

\paragraph{Historical context.}
At the per-element level, the Bernoulli set model reduces to a collection of independent binary channels~\cite{shannonBSC} (symmetric when $\fprate = \fnrate$, asymmetric otherwise), to the misclassification model of Bross~\cite{bross1954misclassification}, and to the randomized-response mechanism of Warner~\cite{warner1965} applied to a membership indicator.
Probabilistic sets~\cite{hirota1981probabilistic} and filters with two-sided error~\cite{donnet2006retouched,laufer2011generalized} also predate this work.
The present paper is an axiomatic restatement of these ideas; the distributions and the composition rule it derives are standard, and it serves as the reference point for the compositional algebra developed in~\cite{bernoulliComposition}.

\paragraph{Further extensions.}
The algebraic structure developed here extends to richer mathematical objects.
The algebra of approximate maps (functions $X \to Y$ with Bernoulli errors) is treated in a companion paper~\cite{bernoulliMaps}, which studies the composition of approximate functions and a space-efficient construction; see that paper's October 2026 errata for the scope of its composition results.
A type-theoretic generalization, treating approximate Booleans, maps, and relations as instances of a single algebraic framework over sum, product, and function types, is developed in companion work on approximate algebraic data types.
